\casdocAbsLabel{目的：}
本公开示例以“东亚流行文化语境下佐佐木希影像传播与审美接受研究”为载体，演示中国科学院博士后出站报告在封面、题名页、摘要、关键词、目录、插图清单、附表清单、正文、参考文献、附录与结尾部分上的完整组织方式。

\casdocAbsLabel{方法：}
示例采用“研究背景、资料整理、传播路径、接受机制、结论建议”的五章结构，并用公开文化材料、占位图表和标准化前后置页模拟正式研究报告的基本排版需求。模板层面继续复用 \texttt{bensz-thesis} 的成熟页眉、目录、图表、参考文献与构建脚本，并在包级注册独立的 \texttt{thesis-cas-postdoc} 模板、profile 与 style，确保单项目分发时仍保持完整模板身份。

\casdocAbsLabel{结果：}
最终得到一套可直接编译的中国科学院博士后出站报告项目脚手架。其封面吸收了《博士后研究报告编写规则》中“分类号 / UDC / 密级 / 编号 / 工作完成日期 / 报告提交日期”的核心字段，同时在题名页中补齐博士后姓名、流动站、专业、起止时间和单位信息。正文示例覆盖章节、表格、示意图、附录、个人简历、博士生和博士后期间成果清单、永久通信地址等关键结构。

\casdocAbsLabel{结论：}
该模板证明，在继续复用现有 thesis 公共基础设施的前提下，博士后研究报告也可以以独立 template ID 的方式稳定交付，同时保持单项目与完整仓库两种使用方式的一致性。

\casdocAbsLabel{关键词：}
博士后研究报告；中国科学院；佐佐木希；东亚流行文化；LaTeX 模板
